\scanpage{113}
Soient $X,Y$ deux préschémas. \uline{Applications géométriques} de $E\subset X$ en $Y$: morphismes de $\operatorname{Dis}(\unsure{E;X})=\coprod_{x\in E}\operatorname{Sp}k(x)$ vers $\operatorname{Dis}(Y)$. \inserted{Supposons $E$ loc. constructible} Une application géométrique est dite \uline{constructible} si \unsure{tt} $x\in X$ a un voisinage ouvert \inserted{affine} $U$, et une partition de $U\cap E$ en parties loc. fermées $Z_i$, telle que $f\mid(Z_i)_{\mathrm{red}}$ soit un \uline{morphisme} $Z_i\longrightarrow Y$. \deleted{Il revient au même de dire que}

Ex: une appl. géom. définie par un morphisme. Si on a $E=\bigcup_{i\in I}E_i$ \unsure{famille} \uline{quelconque} de parties loc. constructibles, alors $f:E\mathrel{\overset{\text{appl. géom.}}{\unsure{-|\longrightarrow}}}Y$ est constr. ssi pour \unsure{$\forall$} $i$, $f\mid E_i$ l'est; cela signifie aussi, \inserted{si $\unsure{X|E}$ est quasi-compact et q. sép.} qu'on peut trouver une \uline{partition} globale de $E$ en parties $Z_i$ constructibles loc. fermées, telles que les $f\mid Z_i$ soient des morphismes. / Définition de $\operatorname{App}_{\mathrm{cons}}((X,E),(\unsure{Y|X},F))$, Transitivité.

\unsure{Ctégorie} \inserted{$(\mathrm{Cons\ Aff})$} des modèles constructibles affines: objets les schémas affines (ou si on préfère, quasi-compacts q. sép.) et morphismes les appl. constructibles. Si $E$ est une partie loc. \deleted{\unsure{préschéma, ou bien}} constructible d'un préschéma $X$, on lui associe un foncteur
\[
\operatorname{Cons}(\unsure{X|Z},E):(\mathrm{Cons\ Aff})^\circ\longrightarrow(\mathrm{Ens})
\]
\margin{112}
